\clearpage
\VerdierPhysicalPageAnchor{0210}%
\begingroup\small\setlength{\parskip}{0pt}\setlength{\emergencystretch}{1.5em}
\setlength{\abovedisplayskip}{4pt plus 2pt minus 2pt}%
\setlength{\belowdisplayskip}{4pt plus 2pt minus 2pt}%

\VerdierSemanticAnchor{ch3d-3-1-4-proof}
Seule reste à démontrer la troisième assertion. Les autres assertions sont
des conséquences immédiates de la proposition (2.4.4). L'assertion 2) montre
aussitôt que la dimension cohomologique de $\mathcal A$ est inférieure ou
égale à la borne supérieure des dimensions injectives (resp.\ projectives) des
objets de $\mathcal A$. Supposons que la dimension cohomologique de
$\mathcal A$ soit égale à $n$ et soit :

\VerdierSemanticAnchor{ch3d-3-1-4-proof-resolution}
\[
I^\bullet=\cdots\longrightarrow0\longrightarrow0\longrightarrow I^0
\longrightarrow I^1\longrightarrow I^2\longrightarrow I^3
\longrightarrow\cdots
\]
une résolution de type injectif d'un objet $Y$ de $\mathcal A$. Pour tout
objet $X$ de $\mathcal A$, on a
$\operatorname{Hom}_{K(\mathcal A)}(X[-n-1],I^\bullet)=0$. En particulier,
en prenant $X=\operatorname{Ker}d_{I^\bullet}^{n+1}$, on montre aussitôt que
le morphisme canonique
$\operatorname{Ker}d_{I^\bullet}^{n+1}\longrightarrow I^{n+1}$ se factorise
à travers $I^n$ :

\VerdierSemanticAnchor{ch3d-3-1-4-proof-factorization}
\[
\begin{tikzcd}[column sep=4.8em,row sep=2.8em]
{} & \operatorname{Ker}d_{I^\bullet}^{n+1}
  \arrow[dl,dashed] \arrow[d] & {} \\
I^n \arrow[r] & I^{n+1} \arrow[r] & \cdots
\end{tikzcd}
\]

\VerdierSemanticAnchor{ch3d-3-1-4-proof-truncation}
Par suite, $\operatorname{Ker}d_{I^\bullet}^{n+1}$ et
$\operatorname{Ker}d_{I^\bullet}^{n}$ sont des facteurs directs de $I^n$ et
sont donc injectifs. Le complexe :
\[
\cdots\longrightarrow0\longrightarrow I^0\longrightarrow I^1
\longrightarrow\cdots\longrightarrow I^{n-1}
\longrightarrow\operatorname{Ker}d_{I^\bullet}^{n}
\longrightarrow0\longrightarrow\cdots
\]
est une résolution injective de $Y$ ; ce qui montre que la dimension
injective de $Y$ est inférieure ou égale à $n$. On démontre de même que la
dimension projective de tout objet $X$ de $\mathcal A$ est inférieure ou
égale à $n$.

\VerdierSemanticAnchor{ch3d-3-1-5}
\medskip\noindent\textbf{3.1.5.}\quad
Soient $X$ et $Y$ deux objets de $\mathcal A$. La proposition (3.1.4) montre
que lorsque la catégorie $\mathcal A$ possède suffisamment d'objets injectifs
ou projectifs les groupes $\operatorname{Ext}_{\mathcal A}^{i}(X,Y)$ que
nous avons définis ici coïncident avec ceux qui sont définis classiquement.

\VerdierSemanticAnchor{ch3d-3-2}
\subsection{La construction de Yoneda.}

\VerdierSemanticAnchor{ch3d-3-2-1}
\noindent\textbf{3.2.1.}\quad
Soient $X$ et $Y$ deux objets de $\mathcal A$. Désignons par
$E_{\mathcal A}^{n}(X,Y)$ ($n\geq1$) l'ensemble des suites exactes :

\VerdierSemanticAnchor{ch3d-3-2-1-sequence}
\[
S^n(X,Y)=0\longrightarrow Y\longrightarrow Z_{n-1}
\longrightarrow Z_{n-2}\longrightarrow\cdots\longrightarrow Z_0
\xrightarrow{u_0}X\longrightarrow0.
\]
L'entier $n$ est appelé la longueur de la suite $S^n(X,Y)$. Associons à tout
élément $S^n(X,Y)$ de $E_{\mathcal A}^{n}(X,Y)$ trois complexes :
\endgroup
